% Generated draft: requires mathematical and source review.
\documentclass[11pt]{article}
\usepackage[margin=25mm]{geometry}
\usepackage{fontspec}
\setmainfont{Latin Modern Roman}
\usepackage{amsmath,amssymb,mathtools,unicode-math}
\setmathfont{Latin Modern Math}
\usepackage{enumitem,needspace,xurl,hyperref}
\setlength{\parindent}{0pt}
\setlength{\parskip}{5pt}
\setlength{\emergencystretch}{3em}
\newcommand{\N}{\mathbb{N}}
\newcommand{\Z}{\mathbb{Z}}
\newcommand{\Q}{\mathbb{Q}}
\newcommand{\R}{\mathbb{R}}
\newcommand{\C}{\mathbb{C}}
\newcommand{\F}{\mathbb{F}}
\newcommand{\A}{\mathbb{A}}
\newcommand{\PP}{\mathbb{P}}
\newcommand{\E}{\mathbb{E}}
\newcommand{\eps}{\varepsilon}
\newcommand{\dd}{\,\mathrm{d}}
\begin{document}
\section*{UnsolvedMath ID: 701155}
\textbf{Source problem code:} catalogue entry
\subsection{Definitions and mathematical statement}
An involution in a group \(G\) is an element of order exactly \(2\). Define the property under consideration by
\[
\exists a,b,c\in G:\quad a^2=b^2=c^2=1_G,\qquad a,b,c\ne 1_G,\qquad ab=ba,\qquad G=\langle a,b,c\rangle.
\]
Thus the three involutions generate \(G\) as an abstract group, and two of them commute.

The problem, proposed by Ya. N. Nuzhin in 2002, asks: \emph{Which adjoint Chevalley groups of normal type over \(\Z\) have this property?}

The supplied record does not define ``normal type'' or specify the construction of the integral adjoint Chevalley groups intended. These conventions must be fixed before the requested classification has a fully precise domain; in particular, an identification with all integral points of an adjoint group scheme is not asserted here.

The supplied source reports the following background results, which are retained as attributed claims rather than independently verified assertions:
\begin{itemize}
\item \(\operatorname{SL}_n(\Z)\) has the property for \(n>13\), attributed to M. C. Tamburini and P. Zucca (1997). This is background about special linear groups and does not itself specify the adjoint family in the question.
\item \(\operatorname{PSL}_n(\Z)\) has the property if and only if \(n>4\), attributed to N. Ya. Nuzhin (2008). The source does not explicitly give the initial range of \(n\) for this assertion.
\item \(\operatorname{PSp}_4(\Z)\) does not have the property. The source says this follows from the corresponding fact for \(\operatorname{PSp}_4(3)\), referring to Archive problem 7.30.
\item The group denoted \(G_2\) has the property, attributed to I. A. Timofeenko (2015).
\item The group denoted \(E_l\) has the property, attributed to I. A. Timofeenko (2017).
\end{itemize}
For the last two claims, the notation is copied from the source: the integral base and group form are implicit in context, and the permitted values of \(l\) in \(E_l\) are not supplied. These qualifications remain unresolved; the background claims do not provide a precise list of the cases still being asked about.
\subsection{Short English statement}
Classify the adjoint Chevalley groups of normal type over \(\Z\) that can be generated by three elements of order exactly two, with two of those generators commuting. The source supplies several attributed positive and negative cases, but leaves the integral-group convention, the meaning of normal type, and the ranks covered by \(E_l\) unspecified. The registry labels the problem open; this draft does not independently verify that status or the cited results.
\subsection{Sources}
\begin{itemize}
\item Local UnsolvedMath registry, supplied complete record with permanent ID \(701155\), exact source problem code \texttt{catalogue entry}, title \emph{Kourovka Notebook Problem 15.67}. The record belongs to the collection \emph{Kourovka Notebook (Issues 1--21)}, described as edited by E. I. Khukhro and V. D. Mazurov. Its supplied status is \texttt{open}; its creation and update timestamps are both \texttt{2026-10-16T14:00:00Z}. Upstream revision: \texttt{d371ee917ac9647fbb06cd6fc529deec6f3dda82}.
\item \emph{The Kourovka Notebook (Unsolved Problems in Group Theory)}, problem 15.67, proposed by Ya. N. Nuzhin, 15th issue (2002); the background also identifies the 21st issue edition (2026), \texttt{arXiv:1401.0300v48}. These source details are supplied attribution and have not been independently checked.
\item M. C. Tamburini and P. Zucca, \emph{J. Algebra} \textbf{195}, no. 2 (1997), 650--661. Supplied reference for the assertion about \(\operatorname{SL}_n(\Z)\), \(n>13\); no article title is supplied.
\item N. Ya. Nuzhin, \emph{Vladikavkaz. Mat. Zh.} \textbf{10}, no. 1 (2008), 68--74 (Russian). Supplied reference for the criterion for \(\operatorname{PSL}_n(\Z)\); the author's initials are retained as printed in this citation.
\item \emph{The Kourovka Notebook}, Archive, problem 7.30. Supplied cross-reference for the negative assertion about \(\operatorname{PSp}_4(3)\) and its stated consequence for \(\operatorname{PSp}_4(\Z)\); the archive entry itself is not supplied.
\item I. A. Timofeenko, \emph{J. Siberian Fed. Univ. Ser. Math. Phys.} \textbf{8}, no. 1 (2015), 104--108. Supplied reference for the assertion about \(G_2\); no article title is supplied.
\item I. A. Timofeenko, \emph{Sibrsk. \`Electron. Mat. Izv.} \textbf{14} (2017), 807--820 (Russian). Supplied reference for the assertion about \(E_l\); the journal spelling is retained from the record, and no article title or range of \(l\) is supplied.
\end{itemize}
\end{document}
